\section{Đạo hàm đa thức [Trung bình]}
\subsection{Phân tích \& Thiết kế (M0)}
    \begin{itemize}
    \item \textbf{Input:}
        \begin{itemize}
            \item File \texttt{poly.csv} chứa 1 dòng: \texttt{degree, a0, a1, ..., an}.
            \item File \texttt{x\_values.csv} chứa các giá trị $x$ cần tính đạo hàm (mỗi dòng 1 số).
        \end{itemize}

    \item \textbf{Xử lý:}
        \begin{itemize}
            \item \textbf{Case biên:}
            \begin{itemize}
                \item \texttt{degree = 0} $\rightarrow$ Đa thức hằng $P(x) = a_0$. Đạo hàm $P'(x) = 0$ (bậc 0, hệ số 0).
                \item Giá trị khuyết (\texttt{NA}, \texttt{NaN}) trong file hệ số $\rightarrow$ coi là 0. Trong file $x$ $\rightarrow$ bỏ qua.
            \end{itemize}
            
            \item \textbf{Module 1 (Đọc dữ liệu):}
            \begin{itemize}
                \item Tái sử dụng logic đọc mảng hệ số phân tách bằng dấu phẩy và đọc mảng $x$ từng dòng (đã xử lý missing values).
            \end{itemize}

            \item \textbf{Module 2 (Tính hệ số đạo hàm):}
            \begin{itemize}
                \item Hàm \texttt{computeDerivCoeffs(vector<double> a)}
                \item Áp dụng công thức: \texttt{deriv[i] = (i + 1) * a[i + 1]} với $i$ chạy từ $0$ đến $n-1$.
                \item Nếu mảng \texttt{a} có kích thước $\le 1$ (tương đương \texttt{degree = 0}), trả về mảng con chỉ chứa \texttt{\{0.0\}}.
            \end{itemize}
            
            \item \textbf{Module 3 (Tính giá trị đa thức):}
            \begin{itemize}
                \item Tái sử dụng hàm \texttt{evalHorner} để tính $P'(x)$ tại các điểm $x$ dựa trên mảng hệ số đạo hàm vừa tìm được.
            \end{itemize}

            \item \textbf{Module 4 (Ghi outputs):}
            \begin{itemize}
                \item Hàm \texttt{writeDerivative(path, vector<double> deriv)}: Ghi file \texttt{derivative.csv} định dạng \texttt{degree\_deriv, b0, b1...}
                \item Hàm \texttt{writeOutput(path, vector<double> x\_vals, vector<double> deriv)}: Ghi file \texttt{output.csv} định dạng \texttt{x,P'(x)}.
                \item Tuân thủ quy định dùng \texttt{fixed} và \texttt{setprecision(6)}.
            \end{itemize}
        \end{itemize}

        \item \textbf{Output:}
        \begin{itemize}
            \item File \texttt{derivative.csv}
            \item File \texttt{output.csv}
        \end{itemize}
    \end{itemize}

\subsection{Mã nguồn C++11 (Tách module)}
    \begin{lstlisting}[language=C++, breaklines=true]
    #include <iostream>
    #include <fstream>
    #include <vector>
    #include <string>
    #include <sstream>
    #include <iomanip>

    using namespace std;

    // MODULE 1: Doc da thuc tu file poly.csv
    bool readPoly(string path, int& degree, vector<double>& a) {
        ifstream file(path);
        if (!file.is_open()) {
            cerr << "Khong mo duoc file " << path << "\n";
            return false;
        }

        string line;
        if (getline(file, line)) {
            stringstream ss(line);
            string token;
            
            if (getline(ss, token, ',')) {
                degree = stoi(token);
            }
            
            while (getline(ss, token, ',')) {
                if (token == "NA" || token == "NaN" || 
                    token == "nan" || token.empty()) {
                    a.push_back(0.0);
                } else {
                    a.push_back(stod(token));
                }
            }
        }
        file.close();
        return true;
    }

    // MODULE 2: Doc cac diem x tu x_values.csv
    bool readXValues(string path, vector<double>& x_vals) {
        ifstream file(path);
        if (!file.is_open()) {
            cerr << "Khong mo duoc file " << path << "\n";
            return false;
        }

        string line;
        while (getline(file, line)) {
            if (line.empty() || line == "NA" || line == "NaN" || line == "nan") {
                continue; 
            }
            x_vals.push_back(stod(line));
        }
        file.close();
        return true;
    }

    // MODULE 3: Tinh he so da thuc dao ham
    vector<double> computeDerivCoeffs(const vector<double>& a) {
        vector<double> deriv;
        
        // Case bien: da thuc hang (bac 0) -> dao ham = 0
        if (a.size() <= 1) {
            deriv.push_back(0.0);
            return deriv;
        }
        
        // Cong thuc: b[i] = (i + 1) * a[i + 1]
        for (size_t i = 0; i < a.size() - 1; i++) {
            deriv.push_back((i + 1) * a[i + 1]);
        }
        
        return deriv;
    }

    // MODULE 4: Tinh gia tri da thuc bang Horner
    double evalHorner(const vector<double>& a, double x) {
        if (a.empty()) return 0.0;
        
        int n = a.size() - 1;
        double res = a[n];
        for (int i = n - 1; i >= 0; i--) {
            res = res * x + a[i];
        }
        return res;
    }

    // MODULE 5: Ghi file derivative.csv
    bool writeDerivative(string path, const vector<double>& deriv) {
        ofstream file(path);
        if (!file.is_open()) {
            cerr << "Khong ghi duoc file " << path << "\n";
            return false;
        }

        int degree_deriv = deriv.empty() ? 0 : deriv.size() - 1;
        file << degree_deriv;
        
        file << fixed << setprecision(6);
        for (double coeff : deriv) {
            file << "," << coeff;
        }
        file << "\n";
        
        file.close();
        return true;
    }

    // MODULE 6: Ghi file output.csv
    bool writeOutput(string path, const vector<double>& x_vals, const vector<double>& deriv) {
        ofstream file(path);
        if (!file.is_open()) {
            cerr << "Khong ghi duoc file " << path << "\n";
            return false;
        }

        file << "x,P'(x)\n";
        file << fixed << setprecision(6);
        
        for (double x : x_vals) {
            double px_prime = evalHorner(deriv, x);
            file << x << "," << px_prime << "\n";
        }

        file.close();
        return true;
    }

    // MAIN
    int main() {
        int degree = 0;
        vector<double> a;
        vector<double> x_vals;

        if (!readPoly("poly.csv", degree, a)) return 1;
        if (!readXValues("x_values.csv", x_vals)) return 1;

        vector<double> deriv = computeDerivCoeffs(a);

        if (!writeDerivative("derivative.csv", deriv)) return 1;
        if (!writeOutput("output.csv", x_vals, deriv)) return 1;

        return 0;
    }
    \end{lstlisting}

\subsection{Kế hoạch kiểm thử (Test Cases)}

    \textbf{Test Case 1 (Thông thường - Bậc 3)}
    
    $P(x) = 1 + 2x + 3x^2 + 4x^3 \Rightarrow P'(x) = 2 + 6x + 12x^2$
    
    \textbf{Input (poly.csv):}
    \begin{lstlisting}
    3,1.0,2.0,3.0,4.0\end{lstlisting}
    \textbf{Input (x\_values.csv):}
    \begin{lstlisting}
    1.0
    -1.0\end{lstlisting}
    
    \textbf{Expected Output (derivative.csv):}
    \begin{lstlisting}
    2,2.000000,6.000000,12.000000\end{lstlisting}
    \textbf{Expected Output (output.csv):}
    \begin{lstlisting}[language={}]
    x,P'(x)
    1.000000,20.000000
    -1.000000,8.000000\end{lstlisting}

    \vspace{0.5cm}

    \textbf{Test Case 2 (Case biên -- Đa thức hằng, degree = 0)}
    
    $P(x) = 5.5 \Rightarrow P'(x) = 0$
    
    \textbf{Input (poly.csv):}
    \begin{lstlisting}
    0,5.5\end{lstlisting}
    \textbf{Input (x\_values.csv):}
    \begin{lstlisting}
    10.0\end{lstlisting}
    
    \textbf{Expected Output (derivative.csv):}
    \begin{lstlisting}
    0,0.000000\end{lstlisting}
    \textbf{Expected Output (output.csv):}
    \begin{lstlisting}[language={}]
    x,P'(x)
    10.000000,0.000000\end{lstlisting} 